\chapter{Discussion}\label{ch:discussion}

\section{What the modest improvement means}

Expression features reduced the error of a tissue-only baseline by 2.85\%, consistently in
direction (nine of ten repeats) but small in size and not uniform across tissues, with
Pancreas moving in the opposite direction by nearly three times the pooled gain. The
defensible interpretation is that \textbf{baseline Hallmark pathway activity carries a small
amount of information about intrinsic radiosensitivity beyond tissue of origin --- on this
endpoint, in this cohort and under this model family}. This is not a demonstration that
expression predicts radiosensitivity usefully, and it is not a validated effect: no interval,
no $p$-value, no negative control and no external cohort stand behind it.

The robustness analyses sharpen this reading rather than soften it. The advantage survives a
change of model family only in diminished form, since the random forest gains less (0.00275
against 0.00383). On transformed outcome scales it keeps its direction but nearly disappears
on the neglog scale; a result whose size depends so strongly on the choice of transformation
is small enough to be shaped by it. The one encouraging result, the leave-one-tissue-out
$R^2$ of $+0.031$ for expression alone, comes from an easier comparison on a different design,
and it reverses for two of the nine well-supported tissues (Pancreas $-0.0142$,
Esophagus/Stomach $-0.0030$). Pancreas is unfavourable in both the primary and the
leave-one-tissue-out analyses, so the tissue-level picture, and not only the pooled gain,
must be part of any statement about generalisation.

Three further cautions apply. The comparison is nested (M2 contains M1 as an offset), so the
errors of the two models are strongly correlated and the paired difference, while the right
statistic, is a narrow one. Repeated cross-validation re-splits one fixed set of components
and cannot estimate generalisation to new cohorts. And the baseline itself is weak: M1 is
marginally worse than the global mean on MAE, which indicates that tissue of origin, as
encoded here, contributes little once shrinkage is taken into account.

\section{Why out-of-sample $R^2$ remains low}

Figure~\ref{fig:prediction} shows predictions compressed near the cohort mean against
observations that span the whole interval. Several causes act together and cannot be
separated with these data. The \textbf{label} is a literature median that aggregates
laboratories, dose rates and protocols, and the field's own reporting audit found a third of
papers omitting replicates or dose rate \parencite{oike2020}; CLARIFi supplies one median per
line, so per-line assay variability cannot even be estimated. The \textbf{sample} is small
(237 independent components) and tissue-imbalanced. The \textbf{features} are 50 coarse
pathway summaries of baseline expression, which may simply not encode the determinants of
radiation response; intrinsic radiosensitivity may depend on dynamic, post-irradiation
responses that a baseline profile cannot capture. And the \textbf{design is strict}: grouping
by relatedness and confining all preprocessing to training folds remove the two most common
sources of inflated performance.

This near-zero $R^2$ closely reproduces RadSigBench, which reported mean out-of-sample $R^2$
of 0.01 and 0.00 for seven published signatures and found them no better than random gene
sets \parencite{oconnor2022}. \textbf{Agreement with the strongest prior negative result is a
substantive finding, not a failure}: it indicates that the pipeline measures what it claims
to measure, and it is evidence against the more optimistic claims of the signature
literature.

\section{The value and uncertainty of expression-derived pathways}

The pathway ranking is best read as a description of a weak fitted model. Its magnitudes are
of the same order as the whole predictive gain, its rank stability across fits is modest
(mean Kendall $\tau$ 0.225, mean top-10 Jaccard 0.391), and correlated gene sets share or
exchange importance. The appearance of HALLMARK\_P53\_PATHWAY at the top is biologically
\emph{plausible}, since p53 governs the damage response, cell-cycle arrest and apoptosis, and
its predominantly negative coefficient is consistent with higher p53-pathway activity
accompanying lower survival. Plausibility is, however, precisely the trap: without a
demonstrated prediction gain and without negative controls, a plausible ranking cannot be
distinguished from a coincidence, and the absence of DNA repair, G2M checkpoint and hypoxia
from the shortlist shows that the ranking is not simply recovering textbook radiobiology. The
ranking is therefore reported as capped, exploratory and hypothesis-generating, and no
mechanistic conclusion is drawn from it.

\section{How the original \textit{Moringa} findings motivate further work}

The original study's lasting contribution is a \emph{question} rather than a measurement this
project could use: the post-irradiation proliferative response to a plant extract differed
between cell types, most clearly for the leaf extract (PEF 1.26 in V79 against 0.90 and 0.99
in the two carcinoma lines). The computational work formalises the general version of that
question on a valid endpoint and with proper grouping, and finds the molecular signal weak.
This does not contradict the original observation, because the endpoints differ: a 72-hour
biomass ratio measured with extract present is not intrinsic clonogenic survival. The
original observation also remains confounded, since the only apparently normal line is the
only non-human one. What the project adds is a defensible framework in which the question can
be asked properly, together with an explicit statement of the experiments needed to answer it.

\section{What compound prioritisation adds, and why the candidates remain hypotheses}

The compound module turns an open-ended phytochemical literature into an auditable evidence
trail: compound $\rightarrow$ documented presence in seed or leaf with a named solvent
$\rightarrow$ primary experimental association with a radiation-response process
$\rightarrow$ the genes actually measured. Its practical value lies as much in what it
\emph{excludes} as in what it ranks. It removes benzyl isothiocyanate (a root, not a seed or
leaf, constituent), gallic and caffeic acid (no primary \textit{Moringa} report), a UVB
literature that is not radiation evidence, carrier-confounded nanoparticle studies and a
retracted paper; it places the free flavonol aglycones in the weakest identification band
rather than excluding them, as their tentative dereplication warrants; and it flags the
tumour-protection liability of chlorogenic acid, which a naive \enquote{antioxidants are
radioprotectors} reading would have turned into a selling point.

The candidates remain hypotheses for specific reasons. No compound has been shown to be
present in the extracts that were tested, because those extracts were never analysed.
Moringin, the most \textit{Moringa}-specific candidate, has no radiation study at all. The
compounds that do have radiation evidence were mostly tested at concentrations or doses whose
relevance to an extract is unknown, several at levels well above plausible dietary exposure,
and many abstracts omit the radiation dose or the concentration. An antioxidant is not
thereby a radioprotector; a compound that reduces tumour-cell viability does not thereby
protect normal tissue; and a high priority score reflects the state of the published
evidence, not a predicted experimental outcome.

The link to the machine-learning results is deliberately kept weak, because the evidence for
it is weak. With the gene-set ranking capped at tier~C, the compound prioritisation rests on
external biological evidence, and any overlap between a compound's measured genes and a
ranked gene set is reported as a hypothesis. \textbf{Nothing in this report claims that a
prioritised pathway and a prioritised compound together imply an effect on clonogenic SF2.}

\section{Separating supported results from plausible explanations}

Table~\ref{tab:supported} separates what the saved outputs support from explanations that are
plausible but not established. The left-hand column is the evidential content of this report;
the right-hand column is the agenda for future work.

\begin{table}[tbp]
  \caption{Results supported by the saved outputs, separated from plausible explanations
    that are not established.}\label{tab:supported}
  \small
  \begin{tabular}{@{}L{0.47\linewidth}L{0.47\linewidth}@{}}
    \toprule
    \thead{Supported by the outputs} & \thead{Plausible, not established} \\
    \midrule
    $\Delta$MAE $=0.00383$ (2.85\%); M2 better in nine of ten repeats
      & That expression features would improve prediction in an independent cohort \\
    Out-of-sample $R^2$ near zero for all five models
      & That richer features or more data would raise it materially \\
    M2 lower error in eight of nine well-supported tissues, worse in Pancreas
      & That tissue-specific differences reflect tissue-specific biology \\
    Expression-only transfer to unseen tissues: $R^2=+0.031$, negative for Pancreas and
      Esophagus/Stomach & That expression generalises to tissues absent from training \\
    Advantage keeps its direction on three outcome scales but nearly vanishes on one
      & That the raw-scale magnitude is the \enquote{true} size of the effect \\
    Seven gene sets meet the tier-C criteria; the p53 pathway ranks first
      & That the p53 pathway influences radiosensitivity in these lines \\
    Importance rank stability is modest across fits
      & That a larger sample would stabilise the ranking \\
    Niazimicin is confirmed in seed and leaf DCM fractions
      & That it was present in the original study's extracts \\
    Moringin has no published ionising-radiation study
      & That it would behave like benzyl isothiocyanate under radiation \\
    Chlorogenic acid protected hepatocellular carcinoma cells from radiation in one study
      & That it would protect human tumours during radiotherapy \\
    \bottomrule
  \end{tabular}
\end{table}
